\section{指定时段与指定日期的完整结果}\label{sec:outputs}
以下表格按照题目表1、表2、表3所要求的字段组织。购电量、充放电量、储电量均以kWh计，费用以元计。第一问采用三组购电时段并列、两组充放电时段并列的原表布局；年度结果对每个指定日期分别列示。第三问和4-3同时给出0点原计划与最终调整后购电，充放电为实际执行量。紧急购电按连续发生区间汇总，四个日期并列，不只列每日合计。

\subsection{问题一指定结果}
\begin{table}[htbp]\centering\scriptsize
\caption{问题一：指定时段及全天购电}
\begin{tabular}{lrrrrr}\toprule
时间段 & 购电量 & 时间段 & 购电量 & 时间段 & 购电量\\\midrule
10:00--10:10 & 0.00 & 12:00--12:10 & 480.41 & 14:00--14:10 & 0.00\\
16:00--16:10 & 445.43 & 18:00--18:10 & 531.89 & 20:00--20:10 & 0.00\\
\multicolumn{2}{l}{全天计划购电量} & 59482.70 & \multicolumn{2}{l}{全天实际购电费} & 35126.95\\
\bottomrule\end{tabular}\end{table}

\begin{table}[htbp]\centering\scriptsize
\caption{问题一：四小时充放电与首末储电量}
\begin{tabular}{lrrrrr}\toprule
时间段 & 充电量 & 放电量 & 时间段 & 充电量 & 放电量\\\midrule
00:00--04:00 & 4500.00 & 0.00 & 04:00--08:00 & 833.33 & 6365.84\\
08:00--12:00 & 4787.96 & 1703.00 & 12:00--16:00 & 5286.04 & 91.10\\
16:00--20:00 & 0.00 & 5780.13 & 20:00--24:00 & 5333.33 & 2859.87\\
\multicolumn{2}{l}{0:00储电量} & 6000.00 & \multicolumn{2}{l}{24:00储电量} & 6000.00\\
\bottomrule\end{tabular}\end{table}

\FloatBarrier
\subsection{问题二指定日期结果}
\begin{table}[htbp]\centering\scriptsize
\caption{问题二 2025-03-20：指定时段及全天购电}
\begin{tabular}{lrrrrr}\toprule
时间段 & 购电量 & 时间段 & 购电量 & 时间段 & 购电量\\\midrule
10:00--10:10 & 0.00 & 12:00--12:10 & 626.63 & 14:00--14:10 & 0.00\\
16:00--16:10 & 494.24 & 18:00--18:10 & 706.68 & 20:00--20:10 & 0.00\\
\multicolumn{2}{l}{全天计划购电量} & 65615.94 & \multicolumn{2}{l}{全天实际购电费} & 43269.00\\
\bottomrule\end{tabular}\end{table}

\begin{table}[htbp]\centering\scriptsize
\caption{问题二 2025-03-20：四小时充放电与首末储电量}
\begin{tabular}{lrrrrr}\toprule
时间段 & 充电量 & 放电量 & 时间段 & 充电量 & 放电量\\\midrule
00:00--04:00 & 1666.67 & 0.00 & 04:00--08:00 & 833.33 & 5501.23\\
08:00--12:00 & 5912.26 & 3138.77 & 12:00--16:00 & 5454.26 & 566.88\\
16:00--20:00 & 0.00 & 5671.59 & 20:00--24:00 & 3498.06 & 2968.41\\
\multicolumn{2}{l}{0:00储电量} & 8550.00 & \multicolumn{2}{l}{24:00储电量} & 4348.26\\
\bottomrule\end{tabular}\end{table}

\FloatBarrier
\begin{table}[htbp]\centering\scriptsize
\caption{问题二 2025-06-21：指定时段及全天购电}
\begin{tabular}{lrrrrr}\toprule
时间段 & 购电量 & 时间段 & 购电量 & 时间段 & 购电量\\\midrule
10:00--10:10 & 0.00 & 12:00--12:10 & 0.00 & 14:00--14:10 & 0.00\\
16:00--16:10 & 250.70 & 18:00--18:10 & 0.00 & 20:00--20:10 & 0.00\\
\multicolumn{2}{l}{全天计划购电量} & 46473.41 & \multicolumn{2}{l}{全天实际购电费} & 26609.86\\
\bottomrule\end{tabular}\end{table}

\begin{table}[htbp]\centering\scriptsize
\caption{问题二 2025-06-21：四小时充放电与首末储电量}
\begin{tabular}{lrrrrr}\toprule
时间段 & 充电量 & 放电量 & 时间段 & 充电量 & 放电量\\\midrule
00:00--04:00 & 2869.46 & 0.00 & 04:00--08:00 & 833.33 & 2615.32\\
08:00--12:00 & 5087.51 & 0.00 & 12:00--16:00 & 3470.92 & 0.00\\
16:00--20:00 & 0.00 & 6017.99 & 20:00--24:00 & 8166.67 & 2622.01\\
\multicolumn{2}{l}{0:00储电量} & 2670.80 & \multicolumn{2}{l}{24:00储电量} & 8550.00\\
\bottomrule\end{tabular}\end{table}

\FloatBarrier
\begin{table}[htbp]\centering\scriptsize
\caption{问题二 2025-09-23：指定时段及全天购电}
\begin{tabular}{lrrrrr}\toprule
时间段 & 购电量 & 时间段 & 购电量 & 时间段 & 购电量\\\midrule
10:00--10:10 & 0.00 & 12:00--12:10 & 549.48 & 14:00--14:10 & 0.00\\
16:00--16:10 & 548.33 & 18:00--18:10 & 777.56 & 20:00--20:10 & 0.00\\
\multicolumn{2}{l}{全天计划购电量} & 70464.22 & \multicolumn{2}{l}{全天实际购电费} & 48722.41\\
\bottomrule\end{tabular}\end{table}

\begin{table}[htbp]\centering\scriptsize
\caption{问题二 2025-09-23：四小时充放电与首末储电量}
\begin{tabular}{lrrrrr}\toprule
时间段 & 充电量 & 放电量 & 时间段 & 充电量 & 放电量\\\midrule
00:00--04:00 & 1666.67 & 0.00 & 04:00--08:00 & 833.33 & 5644.61\\
08:00--12:00 & 6191.30 & 2995.39 & 12:00--16:00 & 4841.85 & 296.85\\
16:00--20:00 & 0.00 & 5695.38 & 20:00--24:00 & 8166.67 & 2944.62\\
\multicolumn{2}{l}{0:00储电量} & 8550.00 & \multicolumn{2}{l}{24:00储电量} & 8550.00\\
\bottomrule\end{tabular}\end{table}

\FloatBarrier
\begin{table}[htbp]\centering\scriptsize
\caption{问题二 2025-12-21：指定时段及全天购电}
\begin{tabular}{lrrrrr}\toprule
时间段 & 购电量 & 时间段 & 购电量 & 时间段 & 购电量\\\midrule
10:00--10:10 & 0.00 & 12:00--12:10 & 989.93 & 14:00--14:10 & 0.00\\
16:00--16:10 & 863.82 & 18:00--18:10 & 719.77 & 20:00--20:10 & 0.00\\
\multicolumn{2}{l}{全天计划购电量} & 98935.42 & \multicolumn{2}{l}{全天实际购电费} & 65505.96\\
\bottomrule\end{tabular}\end{table}

\begin{table}[htbp]\centering\scriptsize
\caption{问题二 2025-12-21：四小时充放电与首末储电量}
\begin{tabular}{lrrrrr}\toprule
时间段 & 充电量 & 放电量 & 时间段 & 充电量 & 放电量\\\midrule
00:00--04:00 & 1666.67 & 0.00 & 04:00--08:00 & 1666.67 & 1508.33\\
08:00--12:00 & 5833.33 & 7806.67 & 12:00--16:00 & 8241.15 & 2760.33\\
16:00--20:00 & 0.00 & 5609.57 & 20:00--24:00 & 8166.67 & 3030.43\\
\multicolumn{2}{l}{0:00储电量} & 8550.00 & \multicolumn{2}{l}{24:00储电量} & 8550.00\\
\bottomrule\end{tabular}\end{table}

\FloatBarrier
\begin{table}[htbp]\centering\scriptsize
\caption{问题二：指定日期全部紧急购电区间}
\setlength{\tabcolsep}{3pt}\begin{tabular}{lrlrlrlr}\toprule
\multicolumn{2}{c}{03-20} & \multicolumn{2}{c}{06-21} & \multicolumn{2}{c}{09-23} & \multicolumn{2}{c}{12-21}\\\midrule
时间段 & 电量 & 时间段 & 电量 & 时间段 & 电量 & 时间段 & 电量\\
00:20--00:30 & 12.39 & 00:20--00:30 & 40.44 & 00:30--01:20 & 22.17 & 01:00--01:20 & 29.05\\
01:10--01:40 & 37.48 & 01:20--01:30 & 21.10 & 02:00--02:20 & 55.41 & 02:10--02:20 & 0.61\\
03:00--03:30 & 70.17 & 01:50--02:10 & 13.24 & 03:50--04:00 & 11.95 & 03:00--03:10 & 15.38\\
04:50--05:10 & 48.10 & 02:50--03:30 & 99.52 & 04:50--05:10 & 34.28 & 03:30--03:50 & 108.42\\
05:50--06:00 & 7.85 & 18:50--19:10 & 78.17 & 05:40--05:50 & 10.72 & 05:50--06:00 & 8.04\\
08:50--09:10 & 26.14 & 20:20--20:30 & 15.05 & 06:20--06:30 & 1.78 & 06:10--06:30 & 41.11\\
11:10--11:20 & 8.27 & 22:00--22:30 & 71.19 & 10:00--10:20 & 13.05 & 10:30--11:30 & 271.92\\
12:30--12:50 & 26.69 & — & — & 10:50--11:00 & 5.98 & 15:20--15:30 & 13.09\\
15:30--16:20 & 208.66 & — & — & 14:30--15:20 & 272.38 & 18:40--19:00 & 8.15\\
16:30--17:20 & 61.84 & — & — & 15:50--16:50 & 147.40 & 19:40--20:10 & 63.76\\
18:40--19:10 & 73.67 & — & — & 17:30--17:40 & 16.92 & — & —\\
20:00--20:20 & 41.73 & — & — & 18:10--18:50 & 111.43 & — & —\\
21:10--21:30 & 42.68 & — & — & 19:10--19:30 & 45.49 & — & —\\
23:30--23:50 & 24.14 & — & — & 19:50--20:30 & 140.64 & — & —\\
— & — & — & — & 20:50--21:10 & 106.59 & — & —\\
— & — & — & — & 21:20--22:00 & 58.24 & — & —\\
— & — & — & — & 22:20--22:30 & 57.72 & — & —\\
— & — & — & — & 23:10--23:40 & 67.37 & — & —\\
合计 & 689.80 & 合计 & 338.71 & 合计 & 1179.52 & 合计 & 559.54\\
\bottomrule\end{tabular}\end{table}

\FloatBarrier
\subsection{问题三指定日期结果}
\begin{table}[htbp]\centering\scriptsize
\caption{问题三 2025-03-20：指定时段及全天购电}
\begin{tabular}{lrrrrr}\toprule
时间段 & 购电量 & 时间段 & 购电量 & 时间段 & 购电量\\\midrule
\multicolumn{6}{c}{0点原计划购电量}\\
10:00--10:10 & 0.00 & 12:00--12:10 & 551.93 & 14:00--14:10 & 0.00\\
16:00--16:10 & 544.34 & 18:00--18:10 & 705.08 & 20:00--20:10 & 0.00\\
\multicolumn{2}{l}{全天计划购电量} & 64668.95 & \multicolumn{2}{l}{全天实际购电费} & 43286.36\\
\multicolumn{6}{c}{最终调整后购电量（不含紧急购电）}\\
10:00--10:10 & 0.00 & 12:00--12:10 & 551.93 & 14:00--14:10 & 0.00\\
16:00--16:10 & 544.34 & 18:00--18:10 & 705.08 & 20:00--20:10 & 0.00\\
\multicolumn{5}{l}{全天最终调整后购电量} & 64101.47\\
\bottomrule\end{tabular}\end{table}

\begin{table}[htbp]\centering\scriptsize
\caption{问题三 2025-03-20：四小时充放电与首末储电量}
\begin{tabular}{lrrrrr}\toprule
时间段 & 充电量 & 放电量 & 时间段 & 充电量 & 放电量\\\midrule
00:00--04:00 & 1629.40 & 0.00 & 04:00--08:00 & 833.33 & 5973.45\\
08:00--12:00 & 6049.15 & 2666.55 & 12:00--16:00 & 5195.34 & 468.04\\
16:00--20:00 & 2.26 & 5675.55 & 20:00--24:00 & 3395.02 & 2974.39\\
\multicolumn{2}{l}{0:00储电量} & 8583.54 & \multicolumn{2}{l}{24:00储电量} & 4246.51\\
\bottomrule\end{tabular}\end{table}

\FloatBarrier
\begin{table}[htbp]\centering\scriptsize
\caption{问题三 2025-06-21：指定时段及全天购电}
\begin{tabular}{lrrrrr}\toprule
时间段 & 购电量 & 时间段 & 购电量 & 时间段 & 购电量\\\midrule
\multicolumn{6}{c}{0点原计划购电量}\\
10:00--10:10 & 0.00 & 12:00--12:10 & 0.00 & 14:00--14:10 & 0.00\\
16:00--16:10 & 223.91 & 18:00--18:10 & 0.00 & 20:00--20:10 & 0.00\\
\multicolumn{2}{l}{全天计划购电量} & 43570.65 & \multicolumn{2}{l}{全天实际购电费} & 25202.45\\
\multicolumn{6}{c}{最终调整后购电量（不含紧急购电）}\\
10:00--10:10 & 0.00 & 12:00--12:10 & 0.00 & 14:00--14:10 & 0.00\\
16:00--16:10 & 176.27 & 18:00--18:10 & 0.00 & 20:00--20:10 & 0.00\\
\multicolumn{5}{l}{全天最终调整后购电量} & 43321.19\\
\bottomrule\end{tabular}\end{table}

\begin{table}[htbp]\centering\scriptsize
\caption{问题三 2025-06-21：四小时充放电与首末储电量}
\begin{tabular}{lrrrrr}\toprule
时间段 & 充电量 & 放电量 & 时间段 & 充电量 & 放电量\\\midrule
00:00--04:00 & 762.83 & 0.00 & 04:00--08:00 & 857.88 & 2271.65\\
08:00--12:00 & 6553.33 & 0.00 & 12:00--16:00 & 4083.07 & 0.00\\
16:00--20:00 & 0.00 & 5995.04 & 20:00--24:00 & 8149.49 & 2645.72\\
\multicolumn{2}{l}{0:00储电量} & 2292.66 & \multicolumn{2}{l}{24:00储电量} & 8533.70\\
\bottomrule\end{tabular}\end{table}

\FloatBarrier
\begin{table}[htbp]\centering\scriptsize
\caption{问题三 2025-09-23：指定时段及全天购电}
\begin{tabular}{lrrrrr}\toprule
时间段 & 购电量 & 时间段 & 购电量 & 时间段 & 购电量\\\midrule
\multicolumn{6}{c}{0点原计划购电量}\\
10:00--10:10 & 0.00 & 12:00--12:10 & 502.56 & 14:00--14:10 & 0.00\\
16:00--16:10 & 560.29 & 18:00--18:10 & 767.73 & 20:00--20:10 & 0.00\\
\multicolumn{2}{l}{全天计划购电量} & 70219.95 & \multicolumn{2}{l}{全天实际购电费} & 47566.49\\
\multicolumn{6}{c}{最终调整后购电量（不含紧急购电）}\\
10:00--10:10 & 0.00 & 12:00--12:10 & 430.27 & 14:00--14:10 & 0.00\\
16:00--16:10 & 560.29 & 18:00--18:10 & 767.73 & 20:00--20:10 & 0.00\\
\multicolumn{5}{l}{全天最终调整后购电量} & 69537.44\\
\bottomrule\end{tabular}\end{table}

\begin{table}[htbp]\centering\scriptsize
\caption{问题三 2025-09-23：四小时充放电与首末储电量}
\begin{tabular}{lrrrrr}\toprule
时间段 & 充电量 & 放电量 & 时间段 & 充电量 & 放电量\\\midrule
00:00--04:00 & 1682.24 & 0.00 & 04:00--08:00 & 833.33 & 5844.03\\
08:00--12:00 & 6049.67 & 2598.82 & 12:00--16:00 & 4754.75 & 331.14\\
16:00--20:00 & 42.16 & 5706.45 & 20:00--24:00 & 8152.71 & 2962.87\\
\multicolumn{2}{l}{0:00储电量} & 8535.99 & \multicolumn{2}{l}{24:00储电量} & 8517.90\\
\bottomrule\end{tabular}\end{table}

\FloatBarrier
\begin{table}[htbp]\centering\scriptsize
\caption{问题三 2025-12-21：指定时段及全天购电}
\begin{tabular}{lrrrrr}\toprule
时间段 & 购电量 & 时间段 & 购电量 & 时间段 & 购电量\\\midrule
\multicolumn{6}{c}{0点原计划购电量}\\
10:00--10:10 & 0.00 & 12:00--12:10 & 965.84 & 14:00--14:10 & 0.00\\
16:00--16:10 & 872.53 & 18:00--18:10 & 708.71 & 20:00--20:10 & 0.00\\
\multicolumn{2}{l}{全天计划购电量} & 98380.06 & \multicolumn{2}{l}{全天实际购电费} & 65090.39\\
\multicolumn{6}{c}{最终调整后购电量（不含紧急购电）}\\
10:00--10:10 & 0.00 & 12:00--12:10 & 954.33 & 14:00--14:10 & 0.00\\
16:00--16:10 & 872.53 & 18:00--18:10 & 708.71 & 20:00--20:10 & 0.00\\
\multicolumn{5}{l}{全天最终调整后购电量} & 98056.59\\
\bottomrule\end{tabular}\end{table}

\begin{table}[htbp]\centering\scriptsize
\caption{问题三 2025-12-21：四小时充放电与首末储电量}
\begin{tabular}{lrrrrr}\toprule
时间段 & 充电量 & 放电量 & 时间段 & 充电量 & 放电量\\\midrule
00:00--04:00 & 1640.14 & 0.00 & 04:00--08:00 & 1660.82 & 1605.72\\
08:00--12:00 & 5884.07 & 7358.10 & 12:00--16:00 & 7766.93 & 2767.14\\
16:00--20:00 & 7.42 & 5602.70 & 20:00--24:00 & 8200.89 & 3043.38\\
\multicolumn{2}{l}{0:00储电量} & 8573.87 & \multicolumn{2}{l}{24:00储电量} & 8576.96\\
\bottomrule\end{tabular}\end{table}

\FloatBarrier
\begin{table}[htbp]\centering\scriptsize
\caption{问题三：指定日期全部紧急购电区间}
\setlength{\tabcolsep}{3pt}\begin{tabular}{lrlrlrlr}\toprule
\multicolumn{2}{c}{03-20} & \multicolumn{2}{c}{06-21} & \multicolumn{2}{c}{09-23} & \multicolumn{2}{c}{12-21}\\\midrule
时间段 & 电量 & 时间段 & 电量 & 时间段 & 电量 & 时间段 & 电量\\
00:20--00:30 & 12.39 & 00:20--00:30 & 40.44 & 00:30--01:20 & 22.17 & 01:00--01:20 & 29.05\\
01:10--01:40 & 37.48 & 01:20--01:30 & 21.10 & 02:00--02:20 & 55.41 & 02:10--02:20 & 0.61\\
03:00--03:30 & 70.17 & 01:50--02:10 & 13.24 & 03:50--04:00 & 11.95 & 03:00--03:10 & 15.38\\
04:50--05:10 & 48.10 & 02:50--03:30 & 99.52 & 04:50--05:10 & 34.31 & 03:30--03:50 & 108.42\\
05:50--06:00 & 6.72 & 18:50--19:10 & 79.29 & 05:40--05:50 & 17.70 & 05:50--06:00 & 8.04\\
08:10--09:20 & 147.65 & 20:20--20:30 & 18.51 & 06:10--06:30 & 0.57 & 06:20--06:30 & 15.63\\
09:40--10:00 & 56.01 & 22:00--22:30 & 70.63 & 07:40--07:50 & 5.03 & 10:30--11:40 & 363.36\\
11:00--11:30 & 114.30 & — & — & 14:20--15:00 & 75.47 & 13:00--13:10 & 1.78\\
11:40--12:00 & 35.27 & — & — & 15:50--16:40 & 126.89 & 18:40--19:00 & 6.49\\
12:30--12:50 & 109.16 & — & — & 17:30--17:50 & 19.82 & 19:40--20:00 & 59.05\\
13:10--13:20 & 0.81 & — & — & 18:10--18:50 & 96.87 & — & —\\
14:00--14:10 & 6.50 & — & — & 19:10--19:30 & 41.82 & — & —\\
15:50--16:10 & 22.60 & — & — & 19:50--20:30 & 131.14 & — & —\\
16:50--17:00 & 0.11 & — & — & 20:50--21:10 & 106.59 & — & —\\
18:40--19:10 & 85.04 & — & — & 21:20--22:00 & 78.37 & — & —\\
19:50--20:20 & 49.50 & — & — & 22:20--22:30 & 57.72 & — & —\\
21:10--21:30 & 47.58 & — & — & 23:10--23:40 & 67.37 & — & —\\
23:30--23:50 & 24.14 & — & — & — & — & — & —\\
合计 & 873.52 & 合计 & 342.74 & 合计 & 949.18 & 合计 & 607.81\\
\bottomrule\end{tabular}\end{table}

\FloatBarrier
\subsection{4-2指定日期结果}
\begin{table}[htbp]\centering\scriptsize
\caption{4-2 2025-03-20：指定时段及全天购电}
\begin{tabular}{lrrrrr}\toprule
时间段 & 购电量 & 时间段 & 购电量 & 时间段 & 购电量\\\midrule
10:00--10:10 & 0.00 & 12:00--12:10 & 653.15 & 14:00--14:10 & 12.32\\
16:00--16:10 & 494.35 & 18:00--18:10 & 0.00 & 20:00--20:10 & 731.74\\
\multicolumn{2}{l}{全天计划购电量} & 63984.26 & \multicolumn{2}{l}{全天实际购电费} & 43741.78\\
\bottomrule\end{tabular}\end{table}

\begin{table}[htbp]\centering\scriptsize
\caption{4-2 2025-03-20：四小时充放电与首末储电量}
\begin{tabular}{lrrrrr}\toprule
时间段 & 充电量 & 放电量 & 时间段 & 充电量 & 放电量\\\midrule
00:00--04:00 & 2500.00 & 0.00 & 04:00--08:00 & 1500.00 & 5904.30\\
08:00--12:00 & 5922.21 & 3144.53 & 12:00--16:00 & 5058.38 & 520.46\\
16:00--20:00 & 0.00 & 6405.44 & 20:00--24:00 & 833.33 & 2234.56\\
\multicolumn{2}{l}{0:00储电量} & 7950.00 & \multicolumn{2}{l}{24:00储电量} & 1950.00\\
\bottomrule\end{tabular}\end{table}

\FloatBarrier
\begin{table}[htbp]\centering\scriptsize
\caption{4-2 2025-06-21：指定时段及全天购电}
\begin{tabular}{lrrrrr}\toprule
时间段 & 购电量 & 时间段 & 购电量 & 时间段 & 购电量\\\midrule
10:00--10:10 & 0.00 & 12:00--12:10 & 0.00 & 14:00--14:10 & 0.00\\
16:00--16:10 & 252.87 & 18:00--18:10 & 487.95 & 20:00--20:10 & 0.00\\
\multicolumn{2}{l}{全天计划购电量} & 49515.72 & \multicolumn{2}{l}{全天实际购电费} & 24101.55\\
\bottomrule\end{tabular}\end{table}

\begin{table}[htbp]\centering\scriptsize
\caption{4-2 2025-06-21：四小时充放电与首末储电量}
\begin{tabular}{lrrrrr}\toprule
时间段 & 充电量 & 放电量 & 时间段 & 充电量 & 放电量\\\midrule
00:00--04:00 & 0.00 & 0.00 & 04:00--08:00 & 2749.81 & 2227.34\\
08:00--12:00 & 7535.00 & 0.00 & 12:00--16:00 & 3131.66 & 0.00\\
16:00--20:00 & 0.00 & 6017.99 & 20:00--24:00 & 9166.67 & 2757.01\\
\multicolumn{2}{l}{0:00储电量} & 1200.00 & \multicolumn{2}{l}{24:00储电量} & 9300.00\\
\bottomrule\end{tabular}\end{table}

\FloatBarrier
\begin{table}[htbp]\centering\scriptsize
\caption{4-2 2025-09-23：指定时段及全天购电}
\begin{tabular}{lrrrrr}\toprule
时间段 & 购电量 & 时间段 & 购电量 & 时间段 & 购电量\\\midrule
10:00--10:10 & 0.00 & 12:00--12:10 & 568.09 & 14:00--14:10 & 0.00\\
16:00--16:10 & 550.06 & 18:00--18:10 & 0.00 & 20:00--20:10 & 0.00\\
\multicolumn{2}{l}{全天计划购电量} & 70272.43 & \multicolumn{2}{l}{全天实际购电费} & 50538.08\\
\bottomrule\end{tabular}\end{table}

\begin{table}[htbp]\centering\scriptsize
\caption{4-2 2025-09-23：四小时充放电与首末储电量}
\begin{tabular}{lrrrrr}\toprule
时间段 & 充电量 & 放电量 & 时间段 & 充电量 & 放电量\\\midrule
00:00--04:00 & 4166.67 & 0.00 & 04:00--08:00 & 833.33 & 5614.15\\
08:00--12:00 & 6150.98 & 3025.85 & 12:00--16:00 & 4882.17 & 296.85\\
16:00--20:00 & 0.00 & 6245.22 & 20:00--24:00 & 4833.33 & 2394.78\\
\multicolumn{2}{l}{0:00储电量} & 6300.00 & \multicolumn{2}{l}{24:00储电量} & 5550.00\\
\bottomrule\end{tabular}\end{table}

\FloatBarrier
\begin{table}[htbp]\centering\scriptsize
\caption{4-2 2025-12-21：指定时段及全天购电}
\begin{tabular}{lrrrrr}\toprule
时间段 & 购电量 & 时间段 & 购电量 & 时间段 & 购电量\\\midrule
10:00--10:10 & 0.00 & 12:00--12:10 & 993.19 & 14:00--14:10 & 0.00\\
16:00--16:10 & 863.82 & 18:00--18:10 & 719.77 & 20:00--20:10 & 0.00\\
\multicolumn{2}{l}{全天计划购电量} & 97387.54 & \multicolumn{2}{l}{全天实际购电费} & 76456.54\\
\bottomrule\end{tabular}\end{table}

\begin{table}[htbp]\centering\scriptsize
\caption{4-2 2025-12-21：四小时充放电与首末储电量}
\begin{tabular}{lrrrrr}\toprule
时间段 & 充电量 & 放电量 & 时间段 & 充电量 & 放电量\\\midrule
00:00--04:00 & 666.67 & 0.00 & 04:00--08:00 & 0.00 & 833.33\\
08:00--12:00 & 5833.33 & 7806.67 & 12:00--16:00 & 8256.22 & 2772.54\\
16:00--20:00 & 0.00 & 5609.30 & 20:00--24:00 & 8333.33 & 3030.70\\
\multicolumn{2}{l}{0:00储电量} & 10200.00 & \multicolumn{2}{l}{24:00储电量} & 8700.00\\
\bottomrule\end{tabular}\end{table}

\FloatBarrier
\begin{table}[htbp]\centering\scriptsize
\caption{4-2：指定日期全部紧急购电区间}
\setlength{\tabcolsep}{3pt}\begin{tabular}{lrlrlrlr}\toprule
\multicolumn{2}{c}{03-20} & \multicolumn{2}{c}{06-21} & \multicolumn{2}{c}{09-23} & \multicolumn{2}{c}{12-21}\\\midrule
时间段 & 电量 & 时间段 & 电量 & 时间段 & 电量 & 时间段 & 电量\\
00:20--00:30 & 10.64 & 00:20--00:30 & 36.51 & 00:30--00:50 & 10.00 & 01:00--01:20 & 26.83\\
01:10--01:40 & 32.72 & 01:20--01:30 & 16.72 & 01:00--01:10 & 5.36 & 03:00--03:10 & 13.56\\
03:00--03:30 & 67.32 & 02:00--02:10 & 10.40 & 02:00--02:20 & 52.79 & 03:30--03:50 & 103.82\\
04:50--05:10 & 47.52 & 02:50--03:30 & 90.63 & 03:50--04:00 & 9.70 & 05:50--06:00 & 4.79\\
05:50--06:00 & 6.72 & 18:50--19:10 & 78.17 & 04:50--05:10 & 31.67 & 06:10--06:30 & 34.39\\
09:00--09:10 & 17.38 & 20:20--20:30 & 15.05 & 05:40--05:50 & 7.01 & 10:30--11:30 & 259.80\\
12:30--12:50 & 3.28 & 22:00--22:30 & 63.81 & 10:00--10:20 & 4.62 & 15:20--15:30 & 11.52\\
15:30--16:20 & 199.58 & — & — & 14:30--15:20 & 260.90 & 18:40--19:00 & 8.15\\
16:30--17:20 & 61.32 & — & — & 15:50--16:50 & 122.31 & 19:40--20:10 & 63.76\\
18:40--19:10 & 78.09 & — & — & 17:30--17:40 & 16.92 & — & —\\
20:00--20:20 & 43.60 & — & — & 18:10--18:50 & 96.79 & — & —\\
21:10--21:30 & 40.07 & — & — & 19:10--19:30 & 41.85 & — & —\\
23:30--23:50 & 23.20 & — & — & 19:50--20:30 & 140.17 & — & —\\
— & — & — & — & 20:50--21:10 & 105.01 & — & —\\
— & — & — & — & 21:20--22:00 & 54.90 & — & —\\
— & — & — & — & 22:20--22:30 & 57.72 & — & —\\
— & — & — & — & 23:10--23:40 & 46.50 & — & —\\
合计 & 631.45 & 合计 & 311.29 & 合计 & 1064.22 & 合计 & 526.64\\
\bottomrule\end{tabular}\end{table}

\FloatBarrier
\subsection{4-3指定日期结果}
\begin{table}[htbp]\centering\scriptsize
\caption{4-3 2025-03-20：指定时段及全天购电}
\begin{tabular}{lrrrrr}\toprule
时间段 & 购电量 & 时间段 & 购电量 & 时间段 & 购电量\\\midrule
\multicolumn{6}{c}{0点原计划购电量}\\
10:00--10:10 & 0.00 & 12:00--12:10 & 545.41 & 14:00--14:10 & 0.00\\
16:00--16:10 & 525.49 & 18:00--18:10 & -0.00 & 20:00--20:10 & 719.50\\
\multicolumn{2}{l}{全天计划购电量} & 67104.27 & \multicolumn{2}{l}{全天实际购电费} & 49343.99\\
\multicolumn{6}{c}{最终调整后购电量（不含紧急购电）}\\
10:00--10:10 & 0.00 & 12:00--12:10 & 545.41 & 14:00--14:10 & 0.00\\
16:00--16:10 & 525.49 & 18:00--18:10 & 0.00 & 20:00--20:10 & 719.50\\
\multicolumn{5}{l}{全天最终调整后购电量} & 66322.18\\
\bottomrule\end{tabular}\end{table}

\begin{table}[htbp]\centering\scriptsize
\caption{4-3 2025-03-20：四小时充放电与首末储电量}
\begin{tabular}{lrrrrr}\toprule
时间段 & 充电量 & 放电量 & 时间段 & 充电量 & 放电量\\\midrule
00:00--04:00 & 7500.00 & 0.00 & 04:00--08:00 & 3166.67 & 6274.00\\
08:00--12:00 & 5993.55 & 2366.00 & 12:00--16:00 & 5097.01 & 343.95\\
16:00--20:00 & 0.73 & 7170.36 & 20:00--24:00 & 65.10 & 1522.37\\
\multicolumn{2}{l}{0:00储电量} & 1200.00 & \multicolumn{2}{l}{24:00储电量} & 1200.00\\
\bottomrule\end{tabular}\end{table}

\FloatBarrier
\begin{table}[htbp]\centering\scriptsize
\caption{4-3 2025-06-21：指定时段及全天购电}
\begin{tabular}{lrrrrr}\toprule
时间段 & 购电量 & 时间段 & 购电量 & 时间段 & 购电量\\\midrule
\multicolumn{6}{c}{0点原计划购电量}\\
10:00--10:10 & 0.00 & 12:00--12:10 & 0.00 & 14:00--14:10 & 0.00\\
16:00--16:10 & 192.84 & 18:00--18:10 & 453.05 & 20:00--20:10 & 0.00\\
\multicolumn{2}{l}{全天计划购电量} & 34666.48 & \multicolumn{2}{l}{全天实际购电费} & 19165.30\\
\multicolumn{6}{c}{最终调整后购电量（不含紧急购电）}\\
10:00--10:10 & 0.00 & 12:00--12:10 & 0.00 & 14:00--14:10 & 0.00\\
16:00--16:10 & 134.07 & 18:00--18:10 & 453.05 & 20:00--20:10 & 0.00\\
\multicolumn{5}{l}{全天最终调整后购电量} & 34353.81\\
\bottomrule\end{tabular}\end{table}

\begin{table}[htbp]\centering\scriptsize
\caption{4-3 2025-06-21：四小时充放电与首末储电量}
\begin{tabular}{lrrrrr}\toprule
时间段 & 充电量 & 放电量 & 时间段 & 充电量 & 放电量\\\midrule
00:00--04:00 & 0.00 & 0.00 & 04:00--08:00 & 2586.28 & 2075.00\\
08:00--12:00 & 6553.33 & 0.00 & 12:00--16:00 & 4107.81 & 15.41\\
16:00--20:00 & 0.00 & 6096.96 & 20:00--24:00 & 53.86 & 2586.67\\
\multicolumn{2}{l}{0:00储电量} & 1200.00 & \multicolumn{2}{l}{24:00储电量} & 1200.00\\
\bottomrule\end{tabular}\end{table}

\FloatBarrier
\begin{table}[htbp]\centering\scriptsize
\caption{4-3 2025-09-23：指定时段及全天购电}
\begin{tabular}{lrrrrr}\toprule
时间段 & 购电量 & 时间段 & 购电量 & 时间段 & 购电量\\\midrule
\multicolumn{6}{c}{0点原计划购电量}\\
10:00--10:10 & 0.00 & 12:00--12:10 & 456.65 & 14:00--14:10 & 0.00\\
16:00--16:10 & 536.21 & 18:00--18:10 & 0.00 & 20:00--20:10 & 0.00\\
\multicolumn{2}{l}{全天计划购电量} & 66826.12 & \multicolumn{2}{l}{全天实际购电费} & 52706.02\\
\multicolumn{6}{c}{最终调整后购电量（不含紧急购电）}\\
10:00--10:10 & 0.00 & 12:00--12:10 & 306.14 & 14:00--14:10 & 0.00\\
16:00--16:10 & 536.21 & 18:00--18:10 & 0.00 & 20:00--20:10 & 0.00\\
\multicolumn{5}{l}{全天最终调整后购电量} & 66054.20\\
\bottomrule\end{tabular}\end{table}

\begin{table}[htbp]\centering\scriptsize
\caption{4-3 2025-09-23：四小时充放电与首末储电量}
\begin{tabular}{lrrrrr}\toprule
时间段 & 充电量 & 放电量 & 时间段 & 充电量 & 放电量\\\midrule
00:00--04:00 & 7333.33 & 0.00 & 04:00--08:00 & 3333.33 & 6441.96\\
08:00--12:00 & 6025.31 & 2198.04 & 12:00--16:00 & 4936.43 & 245.51\\
16:00--20:00 & 25.85 & 6307.45 & 20:00--24:00 & 16.57 & 2360.41\\
\multicolumn{2}{l}{0:00储电量} & 1200.00 & \multicolumn{2}{l}{24:00储电量} & 1200.00\\
\bottomrule\end{tabular}\end{table}

\FloatBarrier
\begin{table}[htbp]\centering\scriptsize
\caption{4-3 2025-12-21：指定时段及全天购电}
\begin{tabular}{lrrrrr}\toprule
时间段 & 购电量 & 时间段 & 购电量 & 时间段 & 购电量\\\midrule
\multicolumn{6}{c}{0点原计划购电量}\\
10:00--10:10 & 0.00 & 12:00--12:10 & 918.26 & 14:00--14:10 & 0.00\\
16:00--16:10 & 843.15 & 18:00--18:10 & 703.41 & 20:00--20:10 & 0.00\\
\multicolumn{2}{l}{全天计划购电量} & 95873.43 & \multicolumn{2}{l}{全天实际购电费} & 77661.40\\
\multicolumn{6}{c}{最终调整后购电量（不含紧急购电）}\\
10:00--10:10 & 0.00 & 12:00--12:10 & 918.26 & 14:00--14:10 & 0.00\\
16:00--16:10 & 844.01 & 18:00--18:10 & 703.41 & 20:00--20:10 & 0.00\\
\multicolumn{5}{l}{全天最终调整后购电量} & 95320.49\\
\bottomrule\end{tabular}\end{table}

\begin{table}[htbp]\centering\scriptsize
\caption{4-3 2025-12-21：四小时充放电与首末储电量}
\begin{tabular}{lrrrrr}\toprule
时间段 & 充电量 & 放电量 & 时间段 & 充电量 & 放电量\\\midrule
00:00--04:00 & 10666.67 & 0.00 & 04:00--08:00 & 3.60 & 864.86\\
08:00--12:00 & 5833.33 & 7351.09 & 12:00--16:00 & 7885.43 & 2899.17\\
16:00--20:00 & 12.56 & 5655.47 & 20:00--24:00 & 44.62 & 3030.85\\
\multicolumn{2}{l}{0:00储电量} & 1200.00 & \multicolumn{2}{l}{24:00储电量} & 1200.00\\
\bottomrule\end{tabular}\end{table}

\FloatBarrier
\begin{table}[htbp]\centering\scriptsize
\caption{4-3：指定日期全部紧急购电区间}
\setlength{\tabcolsep}{3pt}\begin{tabular}{lrlrlrlr}\toprule
\multicolumn{2}{c}{03-20} & \multicolumn{2}{c}{06-21} & \multicolumn{2}{c}{09-23} & \multicolumn{2}{c}{12-21}\\\midrule
时间段 & 电量 & 时间段 & 电量 & 时间段 & 电量 & 时间段 & 电量\\
00:20--00:30 & 22.46 & 00:10--00:30 & 74.72 & 00:30--01:20 & 111.02 & 00:00--00:20 & 7.92\\
01:00--02:00 & 82.55 & 01:20--02:20 & 123.00 & 01:40--02:30 & 112.37 & 01:00--01:20 & 57.62\\
02:30--02:40 & 0.58 & 02:50--03:30 & 149.65 & 03:40--04:00 & 34.84 & 01:30--01:40 & 11.18\\
03:00--03:40 & 111.77 & 03:50--04:10 & 13.90 & 04:50--05:20 & 67.28 & 02:10--02:20 & 15.66\\
04:50--05:20 & 81.18 & 17:20--17:40 & 14.78 & 05:30--05:50 & 29.41 & 02:40--02:50 & 0.88\\
05:50--06:00 & 20.46 & 18:50--19:10 & 105.08 & 06:10--06:30 & 28.01 & 03:00--03:20 & 37.48\\
06:40--07:00 & 7.10 & 20:20--20:40 & 52.87 & 06:50--07:10 & 12.98 & 03:30--03:50 & 144.79\\
08:00--09:20 & 342.10 & 21:40--22:40 & 125.64 & 07:40--08:00 & 26.75 & 04:00--04:20 & 11.40\\
09:40--10:00 & 80.43 & — & — & 09:30--09:50 & 23.79 & 05:40--06:30 & 97.47\\
10:20--10:30 & 17.89 & — & — & 10:00--10:20 & 7.35 & 07:10--07:20 & 3.74\\
10:50--12:00 & 461.66 & — & — & 10:40--11:00 & 35.51 & 07:50--08:00 & 2.03\\
12:30--13:30 & 191.60 & — & — & 11:30--11:50 & 22.71 & 08:30--09:00 & 24.22\\
14:00--14:20 & 35.44 & — & — & 12:50--13:00 & 0.16 & 09:30--09:50 & 39.29\\
14:50--15:00 & 7.47 & — & — & 13:50--14:00 & 14.31 & 10:20--11:50 & 574.75\\
15:30--16:20 & 90.12 & — & — & 14:20--15:20 & 217.15 & 12:00--12:20 & 39.04\\
16:40--17:00 & 27.47 & — & — & 15:50--16:50 & 263.19 & 12:40--13:10 & 63.72\\
18:40--19:20 & 127.95 & — & — & 17:00--17:50 & 73.03 & 15:20--15:30 & 5.86\\
19:30--20:20 & 93.09 & — & — & 18:00--18:50 & 195.57 & 16:20--16:30 & 11.34\\
21:10--21:40 & 75.12 & — & — & 19:10--19:30 & 82.27 & 18:00--18:20 & 11.27\\
23:30--23:50 & 60.55 & — & — & 19:50--22:00 & 562.42 & 18:40--19:00 & 30.66\\
— & — & — & — & 22:10--22:40 & 77.99 & 19:40--20:10 & 88.22\\
— & — & — & — & 23:10--23:40 & 123.17 & 21:50--22:00 & 2.84\\
— & — & — & — & — & — & 22:50--23:00 & 12.06\\
— & — & — & — & — & — & 23:30--23:40 & 2.50\\
合计 & 1936.99 & 合计 & 659.63 & 合计 & 2121.28 & 合计 & 1295.96\\
\bottomrule\end{tabular}\end{table}

\FloatBarrier

\FloatBarrier
\section{支撑材料与复现说明}
本文主结果对应\texttt{result1.xlsx}、\texttt{result2.xlsx}、\texttt{result3.xlsx}、\texttt{result4-2.xlsx}和\texttt{result4-3.xlsx}，分别位于各问正式结果目录。它们包含题面要求的完整策略，本文没有覆盖或重算这些文件。

论文图片和指定结果表由独立资料准备脚本读取已有结果生成；数据映射与SHA-256记录见本稿目录的\texttt{source\_manifest.json}。该脚本仅读取题目、指定主模型相关文件、新图包和效率补充材料，不调用优化器。LaTeX编译与资产校验过程在本稿说明文件中记录。

主模型的计算可依据各问运行说明复现。第二问正式方案需使用有效残差起点7和56天等权；第三问及4-3的四次更新是预设主方案。效率对照材料单独列示来源，不将其视为本次重新完成的实验。参考资料仅列本文实际使用的题面及内部研究材料，尚未加入未经查阅的外部方法文献。
